%% Copernicus Publications Manuscript Preparation Template for LaTeX Submissions
%% ---------------------------------
%% This template should be used for copernicus.cls
%% The class file and some style files are bundled in the Copernicus Latex Package, which can be downloaded from the different journal webpages.
%% For further assistance please contact Copernicus Publications at: production@copernicus.org
%% https://publications.copernicus.org/for_authors/manuscript_preparation.html
%% Please use the following documentclass and journal abbreviations for preprints and final revised papers.
%% 2-column papers and preprints
\documentclass[gmd, manuscript]{copernicus}
%% Journal abbreviations (please use the same for preprints and final revised papers)
% Advances in Geosciences (adgeo)
% Advances in Radio Science (ars)
% Advances in Science and Research (asr)
% Advances in Statistical Climatology, Meteorology and Oceanography (ascmo)
% Aerosol Research (ar)
% Annales Geophysicae (angeo)
% Archives Animal Breeding (aab)
% Atmospheric Chemistry and Physics (acp)
% Atmospheric Measurement Techniques (amt)
% Biogeosciences (bg)
% Climate of the Past (cp)
% DEUQUA Special Publications (deuquasp)
% Earth Surface Dynamics (esurf)
% Earth System Dynamics (esd)
% Earth System Science Data (essd)
% E&G Quaternary Science Journal (egqsj)
% EGUsphere (egusphere) | This is only for EGUsphere preprints submitted without relation to an EGU journal.
% Earth Observation (eo)
% European Journal of Mineralogy (ejm)
% Geochronology (gchron)
% Geographica Helvetica (gh)
% Geoscience Communication (gc)
% Geoscientific Instrumentation, Methods and Data Systems (gi)
% Geoscientific Model Development (gmd)
% History of Geo- and Space Sciences (hgss)
% Hydrology and Earth System Sciences (hess)
% Journal of Bone and Joint Infection (jbji)
% Journal of Environmentally Compatible Air Transport System (jecats)
% Journal of Micropalaeontology (jm)
% Journal of Sensors and Sensor Systems (jsss)
% Magnetic Resonance (mr)
% Mechanical Sciences (ms)
% Natural Hazards and Earth System Sciences (nhess)
% Nonlinear Processes in Geophysics (npg)
% Ocean Science (os)
% Polarforschung - Journal of the German Society for Polar Research (polf)
% Proceedings of the International Association of Hydrological Sciences (piahs)
% Proceedings of the International Ocean Drilling Programme (piodp)
% Safety of Nuclear Waste Disposal (sand)
% Scientific Drilling (sd)
% SOIL (soil)
% Solid Earth (se)
% State of the Planet (sp)
% The Cryosphere (tc)
% Weather and Climate Dynamics (wcd)
% Web Ecology (we)
% Wind Energy Science (wes)
%% \usepackage commands included in the copernicus.cls:
%\usepackage[german, english]{babel}
%\usepackage{tabularx}
%\usepackage{cancel}
%\usepackage{multirow}
%\usepackage{supertabular}
%\usepackage{algorithmic}
%\usepackage{algorithm}
%\usepackage{amsthm}
%\usepackage{float}
%\usepackage{subfig}
%\usepackage{rotating}
\begin{document}
%% ------------------------------------------------------------------
%% TITLE — proposed main title (GMD expects tool/pipeline name + version)
%% ------------------------------------------------------------------
\title{FAIR or FAST? Interoperability versus training efficiency for machine-learning-ready geospatial data on HPC}
%% Alternative title candidates — pick one, delete the rest:
%% (a) From FAIR archive to fast training: storage-layout trade-offs and an
%%     automated conversion pipeline for ML-ready geospatial data cubes on HPC (v1.0)
%% (b) Serving machine-learning-ready Earth-system data at scale: interoperable
%%     and training-optimised Zarr stores, and the bridge between them (v1.0)
%% (c) TerraNostra-Pipeline v1.0: efficient and FAIR storage of a harmonised
%%     1850--2100 land-use data cube for machine learning on HPC
% \Author[affil]{given_name}{surname}
%% NOTE: PLACEHOLDER AUTHOR ORDER — authors listed ALPHABETICALLY by surname.
%% Final order to be set after the lead-authorship decision (plan §6.1).
%% ORCIDs noted in comments; enter them in the Copernicus submission system.
\Author[1]{Mario}{Acosta}                                        %% ORCID 0000-0001-7054-8168
\Author[1]{Pierre Antoine}{Bretonnière}                          %%
\Author[1]{Marina}{Castaño}                                      %% ORCID 0009-0007-6516-4320
\Author[1]{Miguel}{Castrillo Melguizo}                           %%
\Author[1]{Amanda}{Duarte}                                       %% ORCID 0000-0002-9340-958X
\Author[1]{Ferran}{Mirabent Rubinat}                             %%
\Author[1]{Amirpasha}{Mozaffari}                                 %%  0000-0001-6719-0425
\Author[1]{Oriol}{Tintó Prims}                                   %%
\Author[1]{Joan}{Vedrí}                                   %% ORCID 0009-0008-5890-1051
\affil[1]{Earth Science Department, Barcelona Supercomputing Center, Barcelona, 08034, Spain}
%% The [] brackets identify the author with the corresponding affiliation. 1, 2, 3, etc. should be inserted.
%% If an author is deceased, please add \deceased[$Deceased date if applicable$]{$Author number$} (e.g. \deceased[13 November 2015]{2}) at the end of the affiliations. The author number depends on the placement of the author in the author list, e.g. the third author has number 3.
%% If authors contributed equally, please add \equalcontrib{$Author numbers$} (e.g. \equalcontrib{1,3}) at the end of the affiliations. The author number depends on the placement of the author in the author list, e.g. the third author has number 3.
\runningtitle{FAIR or FAST? Interoperability versus training efficiency for machine-learning-ready geospatial data on HPC}
\runningauthor{Acosta et al.} %% placeholder — matches the alphabetical order above; update once the final first author is decided
\received{}
\pubdiscuss{} %% only important for two-stage journals
\revised{}
\accepted{}
\published{}
%% These dates will be inserted by Copernicus Publications during the typesetting process.
\firstpage{1}
\maketitle


\begin{abstract}
\begin{itemize}
\item Problem: the layout that makes data FAIR is not the layout that trains fast.
\item Two representations: interoperable Zarr store vs.\ training-optimised merged store.
\item Result (measured): 7.6$\times$ per-sample on a 16-year record and 19$\times$ on the
  full historical one; 4.6$\times$ end-to-end on one node; a 10-epoch production training
  cut from ${\sim}55$\,h to 2.3\,h against the five-store baseline.
\item Contribution: an automated bridge, plus the released 1850--2100 land-use cube.
\end{itemize}
\end{abstract}


\copyrightstatement{TEXT} %% This section is optional and can be used for copyright transfers.


%% ==================================================================
%% 1  INTRODUCTION
%% ==================================================================
\introduction  %% \introduction[modified heading if necessary]
\label{sec:intro}

Deep learning on Earth-system data at HPC scale is now routinely limited by data ingestion. A
training sample is a small spatial patch drawn at a random place and year, with every input
variable stacked at that point, and an epoch issues millions of such reads. Parallel
filesystems charge a near-constant metadata cost per file opened, so the per-sample cost is
governed by how many chunk files a sample touches, and the storage layout fixes that count.

Two design goals therefore conflict. A store built for reuse keeps one array per variable with
its own attributes and provenance, so a consumer can open a single field. A store built for
training wants every variable a sample needs inside as few files as possible. The same numbers
in the two arrangements differ by an order of magnitude in file count.

The size of that penalty has not been quantified for a production pipeline. Guidance to prefer
fewer, larger chunks is standard \citep{dask_chunks, abernathey2021}, and the weakness of
parallel filesystems under many small reads is documented \citep{nextplatform_hpcfs,
vef2016gpfs}, but how much a reuse-oriented layout costs a real training run, which layer pays
it, and how it scales with the archive remain open.

We answer this for the AI4Land pipeline \citep{mozaffari2026scalable} on MareNostrum~5
\citep{mn5}. The harmonised cube of Sect.~\ref{sec:substrate} exists as an interoperable store
with one array per variable and a training-optimised store with all variables stacked into two
arrays, both holding numerically equivalent data. We decompose the per-sample cost down to a
single file open and measure both layouts end to end.

Under a common reader the merged layout fetches a sample $7.6\times$ faster over a sixteen-year
record and $19.3\times$ faster over the full historical one; a training run on one node
finishes $4.6\times$ sooner; and ten epochs of the production configuration fell from about
55\,h to 2.3\,h against the five-store layout the project began with. The GPUs, previously idle
more than 90\,\% of the time, now set the pace.

\subsection{The interoperability--efficiency tension for ML-ready geospatial data}
\label{sec:tension}

The two poles are defined by their consumers. One reads little data and much description:
which variables exist, in what units, on what grid, derived how. The other reads one patch of
every variable at one place and year, several million times, and consults the description once.

The conflict is arithmetic. Per-sample fetch time factors into the number of chunk files
opened and the cost per open. The layout fixes the first factor: a per-variable store forces
one read per variable per timestep, and each read touches more than one chunk file whenever
the patch straddles a boundary. For a patch of side $P$ drawn at a uniformly random position
from an array chunked at side $C$ in each spatial dimension, with $P<C$, the expected number
of chunks a read overlaps is
%
\begin{equation}
\label{eq:straddle}
\mathbb{E}[\,n_{\mathrm{chunks}}\,] = \left(1 + \frac{P}{C}\right)^{2},
\end{equation}
%
which gives $2.25$ at $P=256$, $C=512$. In the interoperable layout the product reaches about
93 opens per sample.

The second factor is set by the filesystem and by whether the file's metadata is cached, and it
is close to flat in bytes returned. That flatness makes the first factor decisive: splitting
the same bytes across ten times as many files costs close to ten times as much.

Two properties sharpen the tension as a project grows. The penalty scales with the archive,
since a longer record holds proportionally more chunk files and more array structure for the
reader to carry. And it scales with concurrency: each GPU runs its own reader pool, so a
four-GPU node with twelve readers each puts 48 processes on the same metadata service, shared
with every other job on the machine.

The argument holds for any scheme that cuts arrays into separately addressable objects on a
filesystem that charges per object: netCDF split by variable, tiled rasters with one file per
tile and band, or object stores with a fixed per-request latency. What differs between systems
is the size of the charge; the layout still sets how often it is paid.

\subsection{Related work}
\label{sec:related}

Historically, Earth observation and climatological data have been distributed in GRIB or
NetCDF formats \citep{Gowan2022, Kuluke2021}. Both present structural limitations when
integrated into modern machine learning pipelines. GRIB files read data sequentially, which is
inefficient for repeated selective access, limits parallel I/O and lazy loading, and requires
the metadata to be fully parsed before extracting the data \citep{Gowan2022, Glahn2003}. NetCDF
relies on HDF5 and fixes its chunk layout at file creation. Its file handles are tied to
individual processes, which causes I/O issues when multiple workers access the same dataset,
and it lacks native zero-copy tensor export, so data must pass through CPU RAM first
\citep{Huang2023, Russell2006}. Zarr avoids these constraints, with chunked arrays, flexible
metadata and direct memory sharing for AI workflows \citep{Gowan2022, newman2024zarr}, and is
the format used here.

Zarr adoption in the Earth sciences is framed around ARCO storage: publish once, harmonised
and chunked, so any consumer can pull a sub-region over the network \citep{abernathey2021}.
That literature treats chunk geometry as a parameter tuned for a representative analysis
query, and its guidance is to keep chunks large enough to amortise the per-object overhead
\citep{dask_chunks}. A randomly shuffled training loop is a different consumer, touching every
variable at once and reusing almost nothing between reads. Work on feeding chunked stores
straight into training reports large gains once the chunk shape matches the sampling pattern
\citep{earthmover_dataloader, xarray_gpu_pipeline}, but targets cloud object storage or a
direct path into GPU memory.

On supercomputers the storage layer has been a recognised weak point for deep learning for some
years. These systems were provisioned to stream few large files, while training issues
enormous numbers of small scattered reads, and reports from large machines show the metadata
path saturating well before bandwidth \citep{nextplatform_hpcfs}. The usual remedy adds a
layer, staging the working set onto node-local disk \citep{hvac2022}. That route is closed
here, since the production store exceeds the node-local capacity. We reduce how often the cost
is incurred, by changing the layout.

GPFS optimises for large-file bandwidth and routes metadata operations through a dedicated
service that becomes contended once a system holds millions of small files \citep{weka_gpfs}.
A study of file-create performance on the same architecture finds metadata operations scaling
poorly, with a small fraction of files accounting for most of the total time
\citep{vef2016gpfs}. The read-side measurements of Sect.~\ref{sec:diagnosis} are consistent
with this.

%% ---- REMAINING RELATED-WORK THREADS, NOT YET WRITTEN ----
\begin{itemize}
\item FAIR principles and ML dataset metadata (Croissant).
\item Dataloading formats at scale (WebDataset, TFRecord).
\item Our position: the intersection --- trade-off plus tooling.
\end{itemize}

\subsection{Contributions and paper outline}
\begin{itemize}
\item C1: quantified trade-off, with a reusable diagnostic methodology.
\item C2: an automated bridge between representations.
\item C3: the released, DOI'd 1850--2100 cube.
\end{itemize}

%% ==================================================================
%% 2  THE DATASET SUBSTRATE
%% ==================================================================
\section{The dataset substrate: a harmonised 1850--2100 land-use cube}
\label{sec:substrate}
The dataset we consider in this paper is created from seven separate sources, covering the years 1850 to 2100, and in a 1 km grid. Next, we will detail the processing performed to unify the data from these heterogeneous sources. This is also explained in \textit{citation to ML Earth paper?}

\subsection{Source datasets and provenance}

\textbf{Target Labels:} As a target, we use HILDA+ v1.8 \citep{winkler2020hilda+}, distributed via PANGAEA. This contains global land-use at 1 km resolution for the whole world, for the time period 1960--2020. It also expands back to 1899 using Random Forests. It distinguishes 13 classes, although we narrowed it down to 8 by combining all of the different forest sub-types into a single class. This is done since our predictors do not have enough information to distinguish between these different phenologies.
Version note: the pipeline uses HILDA+ v1.0; v2.0 (2025) was released after the pipeline was finalised, and versions are kept consistent with \citet{mozaffari2026scalable}.

\textbf{Predictor Variables:}
\begin{itemize}
    \item \textbf{LUH2} (historical v2h and future v2f) \citep{hurtt2020harmonization}, obtained through input4MIPs/ESGF. This is a global land-use dataset at $0.25^{\circ} \times 0.25^{\circ}$ ($\sim$31~km) resolution, covering the years 1850--2015. We employ the 12 fractional land-use state variables and two secondary-vegetation parameters (mean age and biomass).
    \item \textbf{Köppen-Geiger} climate classification maps \citep{beck2018present, beck2023high} from GLoH2O
    \item \textbf{ISIMIP2b} population data from the ISIMIP HistSoc protocol \citep{frieler2017assessing}
    \item \textbf{GEBCO 2024} elevation and bathymetry \citep{GEBCO2024}
    \item \textbf{PNNL} \citep{li2024global} high-resolution global derived slope and aspect, and soil (texture, organic carbon)
\end{itemize}


\subsection{Spatiotemporal coverage and grid}
To remain consistent across datasets, the input data was aligned to the 1~km WGS84 grid from the HILDA+ dataset.

\textbf{Köppen-Geiger:} Since climatic zones are only available at 30-year intervals, we needed to perform a yearly interpolation that returned valid classes in between the data available. We achieved this by using \textit{stochastic weighted dithering}. For each pixel, we consider the class it has at available time points $t$ and $t+30$. To interpolate for the years in between, we take the starting class and only switch to the ending class if a noise value is lower than a time weight $w(t)$. This weight is assigned so that the probability of switching classes grows linearly across the 30-year interval. We also consider that the pixel will transition a single time in that whole period.

\textbf{ISIMIP2b:} Population data is at $0.5^{\circ} \times 0.5^{\circ}$ resolution and starts at the year 1861. For the spatial issue, we interpolated using population density. Population counts are converted to density ($people/m^2$), conservatively remapped to the target grid, and multiplied by the target cell areas to recover local counts. In terms of the temporal missing data, we used linear back-casting to get data starting on 1850, using the trends from the initial decade (1861–1871) and enforcing a lower limit of zero.

\textbf{PNNL soil data:} In the original dataset there is six-layer vertical soil data. To simplify it, we computed a 2D non-linear depth-weighted average with a decreasing weighting scheme of $w = [0.40, 0.25, 0.15, 0.10, 0.05, 0.05]$.

\textit{Include text on SSP forward extension? As well as tables for the type/units of the data}

\subsection{Licensing and release}
\begin{itemize}
\item All seven sources redistributable with attribution; no share-alike or NC terms.
\item Derived cube released CC BY 4.0 with a DOI (repository TBC).
\item Per-source terms in Table~\ref{tab:licences}.
\end{itemize}

%t
\begin{table*}[t]
\caption{Upstream sources of the released cube, their role, licence, and
redistribution status after modification.}
\label{tab:licences}
\begin{tabular}{llll}
\tophline
Dataset & Role & Licence & Redistribute \\
\middlehline
HILDA+ v1.0 \citep{winkler2020hilda+}        & Target labels             & CC BY 4.0                  & Yes \\
LUH2 v2h/v2f \citep{hurtt2020harmonization}  & Dynamic input             & CC BY 4.0 (CMIP6 terms)    & Yes \\
K{\"o}ppen--Geiger \citep{beck2018present}   & Dynamic input             & CC BY 4.0                  & Yes \\
ISIMIP HistSoc population \citep{frieler2017assessing} & Dynamic input   & CC BY 4.0 (ISIMIP2/2b)     & Yes \\
GEBCO 2024 \citep{GEBCO2024}                 & Static input (elevation)  & Public domain, attribution & Yes \\
PNNL geophysical \citep{li2024global}        & Static input (slope, aspect) & CC0                     & Yes \\
PNNL soil \citep{li2024global}               & Static input (texture, organic) & CC0                  & Yes \\
\bottomhline
\end{tabular}
\belowtable{CMIP6 terms for LUH2 were relaxed to CC BY 4.0 in October 2022.
CC0 is the PNNL DataHub default.}
\end{table*}

%% ==================================================================
%% 3  THE TWO REPRESENTATIONS FAIR <-> FAST
%% ==================================================================
\section{Two representations of one dataset}
\label{sec:two-reps}

Both representations are built from the same cube and carry numerically equivalent data: the
same samples drawn through the same reader agree to floating-point round-off, and the
categorical fields match exactly. They share the 1\,km global grid of
$18\,000 \times 36\,000$ pixels, $512 \times 512$ spatial chunks (2556 per global field), one
chunk per year along time, and Blosc/LZ4 level 5 with byte shuffle. They differ in how
variables are grouped into arrays, in data types and in metadata strategy, which is the
baseline for the trade-off analysis. Table~\ref{tab:layouts} collects the differences.

The project's first production runs used neither. The inputs were spread over five
per-modality Zarr stores, about 680\,GB and 6.73 million chunk files over 1850--2020, and a
sample opened all five in turn. That layout is the production baseline of
Sect.~\ref{sec:throughput}.

\subsection{The interoperable (FAIR) store}
\label{sec:fair-store}

Each variable is kept in its own Zarr array, so its semantic identity and processing history
live in its own attributes and stay traceable without the original sources or the processing
scripts. Any consumer can verify how a field was produced, and standard libraries open and
plot one field without touching the rest.

The store targets the FAIR principles:
\begin{itemize}
    \item \textbf{Findable}: published as a DOI-registered artifact \cite{doi_ref} and indexed
      in the BSC AI Factory data catalog.
    \item \textbf{Accessible}: the layout supports inspection and manipulation through
      standard libraries such as xarray \citep{xarray} and numpy.
    \item \textbf{Interoperable}: attributes are standardised following CF conventions.
    \item \textbf{Reusable}: the metadata carries the origin and processing history of every
      field, providing the context needed for validation.
\end{itemize}

The file count follows from the arrangement. Each of the 17 time-dependent variables holds
$2556 \times T$ chunk files over $T$ years and each of the 6 static variables holds 2556, so
the sixteen-year window 2000--2015 comes to 710\,571 files and 66\,GB and the full historical
record 1901--2015 to 5\,012\,319 files and 435\,GB. Most of those files are small: ocean chunks
compress to the codec floor, putting the median chunk file at about 4\,kB and the mean at about
0.1\,MB.

\subsection{The training-optimised merged store}
\label{sec:merged-store}

Profiling the pipeline on the interoperable store identified an ingestion bottleneck, diagnosed
in Sect.~\ref{sec:diagnosis}. The training-optimised layout is a single store with two
coordinate-defined groups, \emph{static} and \emph{dynamic}, each holding exactly one
consolidated array. Variable names are promoted to a leading coordinate axis, the spatial and
temporal coordinates are preserved, and per-variable descriptors are dropped.

Only the variable axis changes chunking. The dynamic array is
$(17, T, 18\,000, 36\,000)$ chunked $(17, 1, 512, 512)$, so one chunk file carries all 17
dynamic variables for one year and one spatial tile; the static array is
$(6, 18\,000, 36\,000)$ chunked $(6, 512, 512)$. Spatial and temporal chunking and the
compressed bytes are unchanged. File counts fall from 710\,571 to 43\,459 and from
5\,012\,319 to 296\,503, about sixteen-fold, and the surviving files are larger, at a median of
72\,kB and a mean of 1.4\,MB for the dynamic group.

A single array carries a single data type, so the categorical fields are stored as floating
point and cast back at read time. The read pattern collapses along with the file count. A
sample needs every variable at one place across the years its autoregressive sequence spans:
41 reads in the interoperable layout, and in the merged layout two read calls, one on the
dynamic array covering three years and one on the static array, spanning four single-year
accesses in total. With the straddling factor of Eq.~(\ref{eq:straddle}) the per-sample open
count falls from $41 \times 2.25 \approx 92$ to $4 \times 2.25 = 9$. Counting chunk-file
fetches over one thousand random samples gives 92.7 and 9.05, and the measured straddle
distribution matches the predicted 25\,\%, 50\,\%, 25\,\% split over one, two and four chunks.

\subsection{What each representation gives up}
\label{sec:gives-up}

The interoperable store trades read throughput for per-variable semantics and a self-describing
structure, quantified in Sect.~\ref{sec:throughput}. The merged store trades semantic depth and
general-purpose usability for speed: a consumer can no longer open one variable on its own,
the per-variable attributes are gone, the categorical fields lose their type, and the
conventions that make the archive interoperable no longer hold. It is also specialised to one
access pattern, with chunk geometry matched to $256^2$ patch sampling; long time series at
single grid points, or whole global fields of one variable, read worse from it than from the
interoperable store.

The storage cost runs in the unexpected direction. The merged store occupies 387\,GB against
435\,GB over the full record and 59\,GB against 66\,GB over the short one, since the same
compressed bytes in fewer files spend less on per-file overhead. The real cost is the second
copy: a project that publishes one store and trains on the other holds both, so the footprint
roughly doubles. Sect.~\ref{sec:bridge} covers the conversion that produces it.

%t
\begin{table*}[t]
\caption{The two on-disk representations of the same cube. Both use the same grid,
$512\times512$ spatial chunking, one-year chunking along time and the same compressor, and both
return numerically equivalent samples. They differ in how variables are grouped into arrays,
which sets the chunk-file count and the per-sample open count. $T$ is the number of years in
the record; opens per sample were measured over one thousand random samples.}
\label{tab:layouts}
\begin{tabular}{lll}
\tophline
Property & Interoperable store & Merged store \\
\middlehline
Data arrays                    & 23, one per variable                  & 2, one per group \\
Dynamic array shape            & $(T, 18\,000, 36\,000)$, per variable & $(17, T, 18\,000, 36\,000)$ \\
Static array shape             & $(18\,000, 36\,000)$, per variable    & $(6, 18\,000, 36\,000)$ \\
Dynamic chunk shape            & $(1, 512, 512)$                       & $(17, 1, 512, 512)$ \\
Static chunk shape             & $(512, 512)$                          & $(6, 512, 512)$ \\
Chunk files, 2000--2015        & 710\,571                              & 43\,459 \\
Chunk files, 1901--2015        & 5\,012\,319                           & 296\,503 \\
Chunk file size, median / mean & 4\,kB / 0.1\,MB                       & 72\,kB / 1.4\,MB \\
Size on disk, 2000--2015       & 66\,GB                                & 59\,GB \\
Size on disk, 1901--2015       & 435\,GB                               & 387\,GB \\
Reads per sample               & 41                                    & 2 \\
Chunk files opened per sample  & 92.7                                  & 9.05 \\
Per-variable attributes        & preserved                             & dropped \\
Data types                     & native, per variable                  & one per group; categoricals as float \\
\bottomhline
\end{tabular}
\belowtable{The 17 dynamic variables are the 14 LUH2 fields, population, the
K{\"o}ppen--Geiger class and the HILDA+ labels; the 6 static variables are elevation, slope,
aspect and three soil fields.}
\end{table*}

%% ==================================================================
%% 4  THE EFFICIENCY STUDY
%% ==================================================================
\section{Diagnosing the I/O bottleneck}
\label{sec:diagnosis}

An Nsight Systems trace \citep{nsight} of the production loop on the five-store layout, with
NVTX ranges marking each phase of a training step, shows the GPU track empty across the
recorded window while the data-loading phase of one step runs about 76\,s. Over a longer run
the pattern is periodic: twelve batches under a second each, then one batch of about 80\,s,
repeating. Twelve is the reader count, and the period follows it; with twenty readers the slow
batch lands every twentieth. The whole pool exhausts its queue and refills together. Across a
full run the GPUs are idle more than 90\,\% of the time.

The trace does not say what the readers block on. The candidates sit at different layers: the
worker queue, a lock inside the storage library, decompression, byte volume, and the per-file
charge of the filesystem. Two instruments separate them, one removing the model and one
removing the dataset. Numbers follow in Sect.~\ref{sec:results}.

\subsection{Diagnostic methodology}
\label{sec:methodology}

Both instruments measure under the randomly shuffled patch access the training loop
produces, and both separate cold from warm cache state, which accounts for most of the
variance between otherwise identical runs.

\subsubsection{Standalone dataloader profiler}
\label{sec:profiler}

The first instrument replays the training access pattern with model, optimiser and GPU
removed, drawing the same random place-and-year sequence and timestamping every boundary
inside the fetch. Each measurement runs in a fresh subprocess. It has three phases.

Phase~1 fetches fifty samples serially and attributes time to each input source. On the
interoperable store the mean fetch is about 1.25\,s, of which LUH2 takes about 68\,\%, being
the largest variable group and read at two years per sample. The materialised tensor is
9.5\,MB, so one reader delivers about 8\,MB\,s$^{-1}$.

Phase~2 splits each call into the store read and the in-memory work that follows: fill,
normalise, stack, transpose, cast. The read takes about 92\,\% and the in-memory work about
8\,\%. This capped a proposal to bake the normalisation into a third copy of the dataset;
measured against the unbaked store, the baked variant returned indistinguishable per-sample
times, and the copy was never written.

Phase~3 sweeps the reader count. Throughput rises almost linearly at about 0.8 samples\,s
$^{-1}$ per reader, and at twelve readers it equals the reader count divided by the
single-thread fetch time. That excludes the worker queue and the other coordination layers;
the lever is the cost of one fetch.

\subsubsection{Synthetic filesystem probe}
\label{sec:probe}

The second instrument removes the dataset. It writes $N$ files of size $S$ into a directory on
the same volume, reads them back in written or shuffled order, and times every open. Sweeping
$N$ from 1 to $10^4$ and $S$ from 10\,kB to 100\,MB maps the per-open cost against both, and
holding $N \cdot S$ fixed isolates fragmentation at constant volume. A second configuration
builds two pools in one directory, $10^5$ and $10^6$ files, read randomly in one interleaved
sequence so neither is favoured by cache state, to test whether directory size inflates the
per-open cost. A third reproduces the shape of the real problem: two synthetic stores of about
30\,GB with matched total bytes, one with 23 subdirectories of roughly 13\,000 files whose
sizes are drawn from the measured chunk-size distribution and 40 opens per sample, the other
with one subdirectory of 13\,000 files of about 2.3\,MB and 2 opens per sample, read by twelve
concurrent workers in alternating order.

Two calibrations support the probe. Volume bandwidth was measured with a 1\,MB block-copy
benchmark from one stream and from twelve parallel streams. And when the synthetic file sizes
are drawn from the real chunk-size distribution, the synthetic and real per-open distributions
agree at every percentile.

\subsection{Experimental setup: MareNostrum 5}
\label{sec:setup}

All measurements ran on the MareNostrum~5 accelerated partition: two Intel Xeon Platinum
8460Y+ (80 cores), four H100 64\,GB, 512\,GB DDR5, a 480\,GB node-local NVMe and four NDR200
adapters per node, under SLURM \citep{slurm}. One task per node and one process per GPU gives
twenty cores each; at the production setting of twelve readers per GPU a node puts 48 reader
processes on the storage system.

Data sits on the GPFS scratch volume, shared with every other job on the machine, so both
bandwidth and the metadata service are contended. The pipeline is PyTorch \citep{pytorch} with
Accelerate \citep{accelerate} for DDP, and the stores are opened through xarray \citep{xarray}.
The production configuration is batch 8 per GPU, accumulation 4, twelve readers per GPU,
prefetch 4 and pinned memory. Four stores are compared, the interoperable at 66\,GB and
435\,GB and the merged at 59\,GB and 387\,GB over the same two records, plus the 680\,GB
five-store layout. All cover the historical part of the cube; the released cube runs to 2100,
so a production store holds a longer record than measured here.

Because the machine is shared, every configuration is repeated, the fastest and slowest
repetitions are dropped before averaging, jobs are chained so that consecutive jobs alternate
stores, and no two measurement jobs run concurrently against the same store. Surviving
repetitions agree to within a few per cent.

\subsection{Findings}
\label{sec:findings}

\subsubsection{Metadata cost per file opened}
\label{sec:finding-metadata}

Timing the layers of one chunk read separates the candidates (Table~\ref{tab:per-chunk}). The
cold filesystem read costs 19--20\,ms and barely moves with compressed chunk size: a 25\,kB and
a 122\,kB chunk cost the same. Decompression costs 0.3\,ms, about 1\,\% of the total, and
resolving a chunk index to a path is free. The remaining layer, materialising the labelled
array at about 16\,ms, is taken up in Sect.~\ref{sec:finding-size}.

The synthetic probe identifies the filesystem term as metadata. With the record cached, an open
of any file up to about 100\,kB costs about 0.11\,ms, flat as the directory grows from ten to
ten thousand files. Holding volume at 100\,MB, one file reads in 44\,ms while ten thousand
10\,kB files take 1082\,ms. On a cache miss the same open costs about 20\,ms at the median and
23\,ms at p95, matching the real chunk figure, so a cold open costs roughly $200\times$ a warm
one. Directory size leaves it unchanged: pools of $10^5$ and $10^6$ files, read interleaved,
agree within 8\,\% at every percentile. Cache occupancy decides.

Byte volume is excluded by arithmetic. The volume sustains about 556\,MB\,s$^{-1}$
single-stream and 2510\,MB\,s$^{-1}$ over twelve streams; a sample pulls roughly 3\,MB, which
would take about 5\,ms. The observed fetch is 1.25\,s, some $200\times$ slower.

\subsubsection{Opens per sample as a layout property}
\label{sec:finding-layout}

Given a roughly fixed charge per open, the per-sample cost follows from how many opens the
layout forces, and that count is fixed when the store is written. It is the number of
single-year variable accesses times the straddling factor of Eq.~(\ref{eq:straddle}):
$41 \times 2.25 \approx 92$ for the interoperable store, measured at 92.7, and
$4 \times 2.25 = 9$ for the merged store, measured at 9.05.

The production-shape simulation isolates that difference. With matched bytes, matched
chunk-size distributions and only the per-sample open count differing, twelve concurrent
workers on a cold cache pay about 660\,ms per sample on the fragmented arrangement against
about 90\,ms on the consolidated one, reproduced across two independent measurements. Repeating
at a scale where both file populations stay resident inverts the ordering: the fragmented
arrangement warms from 39\,ms to 8.6\,ms per sample and overtakes the consolidated one, whose
larger files still have to be transferred. The gain therefore comes from cache misses on the
metadata path, which a production-scale store cannot avoid, since a full pass sweeps a working
set far larger than the cache holds.

\subsubsection{Growth of the penalty with archive size}
\label{sec:finding-size}

Per-sample fetch time also grows with the total size of the store, though a sample reads the
same slice either way. Extending the record from sixteen to one hundred and fifteen years costs
the interoperable store a factor of 4.15 per sample and the merged store 1.64. Cache-miss rate
turns out to be the wrong explanation: reading the same fourteen-year window from both stores
gives an identical cold-miss rate (0.994), an identical cold-open cost (25.25\,ms) and a ratio
of 0.99 in the time spent reading chunk data.

The cause sits in the reader. Requesting the same already-cached slice from both stores through
the labelled-array layer, with essentially no I/O, runs $6.87\times$ slower against the longer
record, in proportion to the number of chunks the array holds along time. Slicing through this
layer builds and prunes a description of the deferred computation, and that description
enumerates the array's full chunk structure on every access. The cost is therefore proportional
to the array's total chunk count, and it is paid once per read call.

The layout decides how often. The interoperable store issues 41 reads per sample against the
merged store's two, so the same per-read penalty is levied twenty times more often, and the
merging advantage widens with the archive, from $7.6\times$ over sixteen years to
$19.3\times$ over the full record. Extending a published interoperable store in time slows
every training run that reads it.

This cost lives in the reading library and can also be attacked there, by indexing the chunked
arrays directly. That lever is orthogonal to the layout and is out of scope here; every
measurement in Sect.~\ref{sec:results} uses the same reader on both stores.

\subsubsection{Implications for storing training data on HPC systems}
\label{sec:finding-principle}

The findings compose into a cost model. A sample pays one metadata charge per chunk file
opened, where the file count is the number of single-year variable accesses times the
straddling factor of Eq.~(\ref{eq:straddle}) and the charge is about 20\,ms cold and 0.1\,ms
warm, plus one reader-side charge per read call, proportional to the array's total chunk count.
Both leading terms follow from how the layout groups variables into arrays.

The design rules generalise beyond this dataset, format and filesystem. Group the variables a
sample needs into one array. Choose a spatial chunk side at least twice the patch side, keeping
the straddling factor near its floor. Size chunk files so that the per-open charge amortises,
on the order of a megabyte compressed. Target a single-digit open count per sample and treat it
as the quantity to minimise. And verify that a re-laid-out store returns the same numbers
before believing any speed-up measured on it.

%t
\begin{table}[t]
\caption{Cost of one cold chunk read on the GPFS scratch volume, by layer, from two independent
measurements on production chunks that agreed to within a millisecond per layer. The filesystem
read is nearly flat in compressed chunk size, the signature of a metadata round trip.}
\label{tab:per-chunk}
\begin{tabular}{lr}
\tophline
Layer of a single chunk read & Cold cost (ms) \\
\middlehline
Bytes from disk (metadata lookup plus transfer) & 19--20 \\
Materialising the labelled array                & 16 \\
Decompression (Blosc/LZ4)                       & 0.3 \\
Chunk-index lookup                              & ${\sim}0$ \\
\bottomhline
\end{tabular}
\end{table}

%% ==================================================================
%% 5  THE BRIDGE
%% ==================================================================
\section{An automated bridge between representations}
\label{sec:bridge}
%% THE INTELLECTUAL NOVELTY. Claim depends on the plan §6.2 decision:
%% one-directional (FAIR -> training) vs. bidirectional (round-trip).
\begin{itemize}
\item Beyond ``we made training faster'': users should not have to choose.
\item Publish FAIR; derive the training store locally and reproducibly.
\item Scope: one-directional, with round-trip as demonstrated feature or future work.
\end{itemize}
\subsection{Conversion pipeline design}
\begin{itemize}
\item Read $\rightarrow$ stack $\rightarrow$ re-chunk $\rightarrow$ write; parallelised for HPC.
\item Configuration: chunk shape, dtype policy, variable selection, compression.
\item Manifest (source DOI, version, checksums) makes every derived store traceable.
\item Conversion cost quantified in Sect.~6.4.
\end{itemize}
\subsection{Metadata handling and reconstruction}
\begin{itemize}
\item Categorical maps preserved in the manifest for lossless cast-back.
\item Per-variable attributes kept in a sidecar record, not discarded.
\item If bidirectional: state what round-trips exactly and what is lossy.
\item Croissant: demonstrated lightly; full support is future work.
\end{itemize}
\subsection{Usage workflow}
\begin{itemize}
\item Three steps: fetch the DOI'd store, convert, point the dataloader at it.
\item Worked example with commands, runtime, and disk footprint.
\item Guidance on presets for other sampling patterns.
\end{itemize}

%% ==================================================================
%% 6  RESULTS AND BENCHMARKS
%% ==================================================================
\section{Results: the quantified trade-off}
\label{sec:results}
\begin{itemize}
\item Order: throughput, sensitivity, scale-out, costs.
\item All numbers from the setup of Sect.~\ref{sec:setup}.
\end{itemize}

\subsection{Per-sample and end-to-end throughput}
\label{sec:throughput}

The measurements move from one sample to a full training run. All use the same reader on both
layouts under the setup of Sect.~\ref{sec:setup}.

\subsubsection{Cost of one sample}

Table~\ref{tab:per-sample} gives the median time to fetch one sample with a single reader and
no model, over one hundred samples drawn from the same ten-year window in every store, so that
only the layout and the stored record length vary.

At matched record length the merged layout is $7.6\times$ faster over the sixteen-year record
(174 against 1329\,ms) and $19.3\times$ faster over the full historical one (286 against
5513\,ms). The widening is the size effect of Sect.~\ref{sec:finding-size}: extending the
record costs the interoperable layout $4.15\times$ per sample and the merged layout
$1.64\times$. Against the five-store production layout, the merged store over the full record
is about $28\times$ faster per sample (286 against 7970\,ms).

\subsubsection{Delivery rate of the reader pool}

Table~\ref{tab:delivery} puts the full reader pool back, still without the model. Delivery
scales close to linearly with the reader count on both layouts, so the per-sample gap carries
through to a seven- to eightfold gap in delivery rate.

What matters is where each layout sits against GPU demand. One H100 completes a batch of eight
in about 0.15\,s, consuming about 52 samples\,s$^{-1}$. At twelve readers the interoperable
store delivers 8.7 samples\,s$^{-1}$, roughly six times under that rate, so the GPU starves.
The merged store delivers 70.5, slightly above it. That reversal moves the limit off the data
path.

\subsubsection{End-to-end training on one node}

Table~\ref{tab:end-to-end} reports a controlled comparison with the model in the loop, on one
node and the sixteen-year record, at 200 training, 30 validation and 30 test steps per epoch,
three repetitions per layout with the two interleaved so that neither is warmed by its
predecessor. The merged layout completes the same work $4.63\times$ faster, with the training
and validation phases agreeing to about three per cent.

The shape of the run changes as well. On the interoperable layout the per-batch log keeps the
periodic structure of the original trace: seventeen to twenty slow batches per epoch on an
exact twelve-batch cadence, the largest about 34\,s. On the merged layout the periodic stall
disappears, leaving five, three and one isolated slow batches across the three repetitions, at
no fixed cadence.

\subsubsection{Wall time of a full production run}

The last measurement is a complete production training: ten epochs on the merged store over the
full historical record, one node, production batch size and accumulation. It finished in
8420\,s, or 2.34\,h. The same training on the five per-modality stores took about 55\,h as
recorded by the team before this work, so the layout change accounts for about $23\times$ on
the production workload. That figure compares a measurement against the team's own production
timing; repeating a 55\,h run as a control was beyond the allocation, so the controlled
$4.63\times$ of Table~\ref{tab:end-to-end} remains the conservative number.

Table~\ref{tab:production} shows how time distributes across batches. At twelve readers the
median batch is 154\,ms, the GPU's own compute time for eight samples and therefore the floor,
while the mean is 272\,ms and p95 is 1.3\,s, with 7.7\,\% of batches waiting on a refill.
Twenty readers leave the floor untouched and pull the mean to 200\,ms and p95 to 215\,ms, with
1.8\,\% above one second. The residual stalls are the reader-side size term of
Sect.~\ref{sec:finding-size} showing through at full record length; more readers hide it.

The run converged normally, reaching its best validation loss at epoch eight, and the model
reproduces 87.4\,\% of land-use pixels in the first predicted year of a fourteen-year rollout,
declining to 79.2\,\% by the thirteenth. Since both stores return equivalent samples, the
layout change leaves model quality untouched.

%t
\begin{table}[t]
\caption{Median time to fetch one sample, single reader, no model, over one hundred samples
drawn from the same ten-year window in every store. Only the layout and the stored record
length differ.}
\label{tab:per-sample}
\begin{tabular}{llrr}
\tophline
Layout & Record & Size & Median (ms) \\
\middlehline
Five per-modality stores & 1850--2020 & 680\,GB & 7970 \\
Interoperable            & 2000--2015 &  66\,GB & 1329 \\
Interoperable            & 1901--2015 & 435\,GB & 5513 \\
Merged                   & 2000--2015 &  59\,GB &  174 \\
Merged                   & 1901--2015 & 387\,GB &  286 \\
\bottomhline
\end{tabular}
\end{table}

%t
\begin{table}[t]
\caption{Delivery rate in samples\,s$^{-1}$ with the model out of the loop, 80 batches per
configuration on the sixteen-year record. One H100 consumes about 52 samples\,s$^{-1}$ on this
model at a batch of eight.}
\label{tab:delivery}
\begin{tabular}{rrr}
\tophline
Reader processes & Interoperable & Merged \\
\middlehline
 1 & 0.8 &  5.9 \\
 4 & 3.3 & 25.3 \\
12 & 8.7 & 70.5 \\
\bottomhline
\end{tabular}
\end{table}

%t
\begin{table}[t]
\caption{Controlled end-to-end comparison with the model in the loop, one node, sixteen-year
record, production configuration. Each epoch runs 200 training, 30 validation and 30 test
steps; means over three interleaved repetitions per layout.}
\label{tab:end-to-end}
\begin{tabular}{lrrr}
\tophline
Phase & Interoperable (s) & Merged (s) & Ratio \\
\middlehline
Training   & 292.4 & 62.9 & 4.65 \\
Validation &  61.8 & 13.7 & 4.51 \\
Total      & 354.1 & 76.5 & 4.63 \\
\bottomhline
\end{tabular}
\end{table}

%t
\begin{table}[t]
\caption{Per-batch wall time during the ten-epoch production training on the merged store over
the full historical record, one node, batch 8 per GPU. The median is the GPU compute time and
does not move with the reader count; the extra readers buy the tail.}
\label{tab:production}
\begin{tabular}{rrrrr}
\tophline
Readers & Median (ms) & Mean (ms) & p95 (ms) & Batches $>1$\,s \\
\middlehline
12 & 154 & 272 & 1306 & 7.7\,\% \\
20 & 157 & 200 &  215 & 1.8\,\% \\
\bottomhline
\end{tabular}
\belowtable{At the median the pipeline delivers 51.9 samples\,s$^{-1}$ per GPU; at the mean,
29.4 with twelve readers and 40.0 with twenty.}
\end{table}

\subsection{Sensitivity to chunking strategy, store size, and worker count}
\begin{itemize}
\item Chunking: where mis-chunking erases the gain.
\item Store size: the gain grows with the store.
\item Workers: FAIR store saturates; merged store keeps scaling.
\end{itemize}
\subsection{Scale-out behaviour}
\begin{itemize}
\item Aggregate read pressure on GPFS as node count grows.
\item Does the advantage widen or narrow at scale?
\item Guidance: where each configuration becomes filesystem-limited.
\end{itemize}
\subsection{Cost of interoperability, cost of the bridge}
\begin{itemize}
\item Price of FAIR: slowdown when training directly on the interoperable store.
\item Price of the bridge: one-off conversion cost and the duplicate copy.
\item Summary table: the whole trade-off in one view.
\end{itemize}

%% ==================================================================
%% 7  CONCLUSIONS
%% ==================================================================
\conclusions  %% \conclusions[modified heading if necessary]
\begin{itemize}
\item The trade-off is real and now quantified.
\item Principle: store layout is the largest lever; metadata dominates over bandwidth.
\item The bridge resolves it operationally: publish FAIR, derive fast.
\item The released cube gives the community a substrate to reuse.
\item Future work: round-trip, in-place conversion, full Croissant, other filesystems.
\end{itemize}

%% The following commands are for the statements about the availability of data sets and/or software code corresponding to the manuscript.
\codedataavailability{
%% - Data: the 1850--2100 interoperable cube, CC BY 4.0, DOI (repository TBC).
%% - Code: the pipeline, the conversion (bridge) tooling, and the benchmark
%%   harness (profiler + synthetic probe), with an archived release DOI.
%% - Exact versions used for the results in this paper.
TEXT}

%% No appendix yet. If one is added (e.g. store schemas / chunking configurations),
%% put \appendix before it and keep \noappendix after the last appendix section.
\noappendix

\authorcontribution{TEXT} %% mandatory — settle after the lead-authorship decision
\competinginterests{TEXT} %% mandatory even if no competing interests are present
\begin{acknowledgements}
%% CERISE belongs here only, as funding acknowledgement (grant 101082139).
TEXT
\end{acknowledgements}

%% REFERENCES
%% The reference list is compiled as follows:
%\begin{thebibliography}{}
\bibliographystyle{copernicus}
\bibliography{refs}
%\end{thebibliography}
%% Since the Copernicus LaTeX package includes the BibTeX style file copernicus.bst,
%% authors experienced with BibTeX only have to include the following two lines:
%%
%% \bibliographystyle{copernicus}
%% \bibliography{example.bib}
%%
%% URLs and DOIs can be entered in your BibTeX file as:
%%
%% URL = {http://www.xyz.org/~jones/idx_g.htm}
%% DOI = {10.5194/xyz}
%% LITERATURE CITATIONS
%%
%% command                        & example result
%% \citet{jones90}|               & Jones et al. (1990)
%% \citep{jones90}|               & (Jones et al., 1990)
%% \citep{jones90,jones93}|       & (Jones et al., 1990, 1993)
%% \citep[p.~32]{jones90}|        & (Jones et al., 1990, p.~32)
%% \citep[e.g.,][]{jones90}|      & (e.g., Jones et al., 1990)
%% \citep[e.g.,][p.~32]{jones90}| & (e.g., Jones et al., 1990, p.~32)
%% \citeauthor{jones90}|          & Jones et al.
%% \citeyear{jones90}|            & 1990
%% FIGURES
%% When figures and tables are placed at the end of the MS (article in one-column style), please add \clearpage
%% between bibliography and first table and/or figure as well as between each table and/or figure.
% The figure files should be labelled correctly with Arabic numerals (e.g. fig01.jpg, fig02.png).
%% ONE-COLUMN FIGURES
%%f
%\begin{figure}[t]
%\includegraphics[width=8.3cm]{FILE NAME}
%\caption{TEXT}
%\end{figure}
%
%%% TWO-COLUMN FIGURES
%
%%f
%\begin{figure*}[t]
%\includegraphics[width=12cm]{FILE NAME}
%\caption{TEXT}
%\end{figure*}
%
%
%%% TABLES
%%%
%%% The different columns must be seperated with a & command and should
%%% end with \\ to identify the column brake.
%
%%% ONE-COLUMN TABLE
%
%%t
%\begin{table}[t]
%\caption{TEXT}
%\begin{tabular}{column = lcr}
%\tophline
%
%\middlehline
%
%\bottomhline
%\end{tabular}
%\belowtable{} % Table Footnotes
%\end{table}
%
%%% TWO-COLUMN TABLE
%
%%t
%\begin{table*}[t]
%\caption{TEXT}
%\begin{tabular}{column = lcr}
%\tophline
%
%\middlehline
%
%\bottomhline
%\end{tabular}
%\belowtable{} % Table Footnotes
%\end{table*}
%
%%% LANDSCAPE TABLE
%
%%t
%\begin{sidewaystable*}[t]
%\caption{TEXT}
%\begin{tabular}{column = lcr}
%\tophline
%
%\middlehline
%
%\bottomhline
%\end{tabular}
%\belowtable{} % Table Footnotes
%\end{sidewaystable*}
%
%
%%% MATHEMATICAL EXPRESSIONS
%
%%% All papers typeset by Copernicus Publications follow the math typesetting regulations
%%% given by the IUPAC Green Book (IUPAC: Quantities, Units and Symbols in Physical Chemistry,
%%% 2nd Edn., Blackwell Science, available at: http://old.iupac.org/publications/books/gbook/green_book_2ed.pdf, 1993).
%%%
%%% Physical quantities/variables are typeset in italic font (t for time, T for Temperature)
%%% Indices which are not defined are typeset in italic font (x, y, z, a, b, c)
%%% Items/objects which are defined are typeset in roman font (Car A, Car B)
%%% Descriptions/specifications which are defined by itself are typeset in roman font (abs, rel, ref, tot, net, ice)
%%% Abbreviations from 2 letters are typeset in roman font (RH, LAI)
%%% Vectors are identified in bold italic font using \vec{x}
%%% Matrices are identified in bold roman font
%%% Multiplication signs are typeset using the LaTeX commands \times (for vector products, grids, and exponential notations) or \cdot
%%% The character * should not be applied as mutliplication sign
%
%
%%% EQUATIONS
%
%%% Single-row equation
%
%\begin{equation}
%
%\end{equation}
%
%%% Multiline equation
%
%\begin{align}
%& 3 + 5 = 8\\
%& 3 + 5 = 8\\
%& 3 + 5 = 8
%\end{align}
%
%
%%% MATRICES
%
%\begin{matrix}
%x & y & z\\
%x & y & z\\
%x & y & z\\
%\end{matrix}
%
%
%%% ALGORITHM
%
%\begin{algorithm}
%\caption{...}
%\label{a1}
%\begin{algorithmic}
%...
%\end{algorithmic}
%\end{algorithm}
%
%
%%% CHEMICAL FORMULAS AND REACTIONS
%
%%% For formulas embedded in the text, please use \chem{}
%
%%% The reaction environment creates labels including the letter R, i.e. (R1), (R2), etc.
%
%\begin{reaction}
%%% \rightarrow should be used for normal (one-way) chemical reactions
%%% \rightleftharpoons should be used for equilibria
%%% \leftrightarrow should be used for resonance structures
%\end{reaction}
%
%
%%% PHYSICAL UNITS
%%%
%%% Please use \unit{} and apply the exponential notation (e.g. 20\,\unit{W\,m^{-2}})
\end{document}
